\chapter{From Specifications to Oracles: Requirement-Driven Symbolic Execution}\label{sec:paper1}
Network protocols rely on precise specifications that describe how communicating parties should construct messages, how state is expected to evolve during an interaction, and how error conditions must be handled.
These specifications are expressed in RFC documents, where natural language is used to formulate requirements.
Many of these requirements involve relationships across multiple packets.
For example, a field in one message may depend on a value in a previous message, or a sequence must follow a specific ordering that is valid for the selected key exchange method.

Testing such requirements is challenging. 
Traditional fuzzers concentrate on triggering crashes and often operate without awareness of protocol state.
Standard symbolic execution (SE) tools are powerful for analyzing single inputs, yet they do not directly address multi packet interactions or stateful requirements.
When applied naively to protocol implementations, these tools either fail to explore relevant paths or become overwhelmed by the path explosion caused by large input structures.

The work presented in this chapter introduces a technique that applies SE in a requirement driven manner.
The aim is to systematically search for inputs that violate specific RFC requirements and to produce concrete packet sequences that expose the deviation in the implementation under test.
The objective is not limited to detecting crashes. 
It also focuses on revealing non-conformance flaws that may affect security or interoperability.
\section{Technique}
\subsection{Extracting Specification Requirements}
The first step is to extract protocol requirements from the RFC text and express them as logical predicates over the packets exchanged during a protocol session.
RFCs typically use particular keywords such as "MUST", "MUST NOT", "REQUIRED", "SHALL", and "SHALL NOT" to indicate the strictness of a requirement~\cite{RFC2119}.
These statements can be formulated into predicates that describe how packets must be structured and how they relate to one another.

A simple but representative example appears in the DTLS RFC~\cite[p. 13]{RFC6347}, which specifies the uniqueness of record sequence numbers within a session.
Each DTLS record contains a 48-bit \rfc{sequence\_number} field that uniquely identifies the record within a session.
This field is incremented for every record sent, usually starting from zero. 
The receiver uses sequence numbers to process records in the correct order and to detect duplicates, since any two records with the same \rfc{sequence\_number} represent a replay.
%
\rfcdisplayquote{%
For each received record, the receiver MUST verify that the record contains a sequence number that does not duplicate the sequence number of any other record received during the life of this session.}%
%
This requirement serves as a protection against replayed records (packets).
For a set of records $R$ received during a DTLS session, the requirement can be expressed in predicate logic as:
%
\begin{equation}\label{const:rsnuniqueness}\normalsize
	\forall r, r' \in R :\, r \neq r' \implies 
	r\rfcm{.sequence\_number}\, \neq\, r'\rfcm{.sequence\_number}
\end{equation}
%
This formula captures the essential rule that no two records may share the same sequence number.
Requirements of this kind, which relate fields across multiple packets, appear frequently in network protocol specifications. 
Once formulated, they provide a precise basis for guiding the SE engine toward the regions of the input space where potential violations may occur.

\subsection{Symbolic Execution}
Symbolic execution is applied to each tested implementation through a dedicated test harness that replays a pre-captured sequence of records from a valid DTLS handshake and passes them one at a time to the implementation under test. 
Each record is parsed into a structured representation, which enables SE to access and manipulate selected fields. 
The harness provides a deterministic environment by controlling the order of record processing and by removing nondeterministic behavior such as cookie verification or timeouts.

With the test harness in place, the next step is to select which fields should be made symbolic for checking a requirement.
In our technique, only those fields in a record that are relevant to the requirement being tested are made symbolic, while the remaining fields retain their concrete values from the pre-captured handshake.
This selective approach focuses SE on the parts of the input space that influence the requirement, while avoiding possible path explosion if entire records were made symbolic.

To ensure that SE concentrates on violating cases rather than on the vast space of valid inputs, the harness introduces assumptions that encode the negation of the requirement.
These assumptions restrict SE to inputs under which the requirement may be broken.
The final step is to check whether the implementation uses the invalid input in a forbidden way.
This can be done by checking whether protocol state progresses after reception of invalid records.
We approximate this by successful completion of the protocol handshake. 
A failing assertion is inserted at the point in the implementation code where the handshake finishes successfully.
If SE reaches this assertion, it means that the implementation completed the handshake even though it received an invalid record, which indicates a violation of the requirement.

The predicate in \Cref{const:rsnuniqueness} formulates the DTLS requirement that no two records in a session may share the same sequence number.
To apply our technique to evaluate this requirement for a server, the harness loads the sequence of pre-captured records that a DTLS server would receive during its handshake and marks the client-side \rfc{sequence\_number} fields as symbolic.
All other fields remain concrete.

An assume statement is then introduced that indicates that at least two symbolic sequence numbers are equal.
This directs SE to explore precisely those assignments that violate uniqueness.
To evaluate the effect of this invalid input, an assertion is placed at the code location that indicates successful completion of the handshake. 
If SE reaches this assertion, it means that the implementation completed the handshake even though it received a record with a duplicate sequence number. 
The assertion then fails, and the SE engine reports the concrete assignments that triggered this behavior.

\subsection{Test Case Construction and Validation}
When SE discovers a violating execution path, the engine provides concrete values for all symbolic fields involved.
To construct a test case, these values are substituted back into the structured record representations used by the test harness, while all remaining fields retain their original concrete values from the pre-captured handshake.
The resulting records are then serialized into byte arrays that follow the exact wire format of the protocol.

For the sequence number requirement, this means that SE produces a concrete assignment in which two records share the same sequence number. 
These values replace the symbolic fields in the corresponding records, creating a concrete record sequence that instantiates the violation of the predicate in \Cref{const:rsnuniqueness}.

These concrete record sequences are replayed against the unmodified implementation outside the SE environment.
Validation is done by observing whether the unmodified implementation accepts the sequence and completes the handshake.
Since handshake completion is used as an approximation for protocol progress, a successful handshake in the presence of an invalid record constitutes a confirmed violation of the corresponding requirement.
In the sequence number example, if the server completes its handshake despite the duplicate sequence number, the requirement is violated and the test case is confirmed.

\section{Summary of Results}
The technique was applied to four implementations of the DTLS protocol. These included OpenSSL, MbedTLS, and two variants of TinyDTLS. A collection of sixteen requirements from the DTLS RFC~\cite{RFC6347} were formulated and tested. These requirements ranged from record layer constraints, such as content type, version fields, and epoch handling, to handshake layer rules related to message ordering, fragmentation, and message sequence numbers.

The evaluation revealed a total of thirty six unique bugs, seven of which were classified as vulnerabilities. Several issues were memory related, including out of bounds accesses during message reassembly. Others concerned logical violations of the specification, such as incorrect validation of version fields, improper handling of repeated handshake messages, and failures to respect message ordering rules. Many issues were identified within a short symbolic exploration time, and all of them were acknowledged and fixed by the developers.

In addition to uncovering new issues, our technique also reproduced a known exploitable vulnerability in OpenSSL's function for reassembling fragmented messages.
This reproduction confirms the power and promise of our technique.

For a detailed description of the technique's implementation, the full set of requirements we evaluated, and a comprehensive analysis of the issues we discovered, the reader is referred to \textcolor{red}{Paper 1} included later in this thesis.
\section{Limitations}
The technique presented in this chapter has two main limitations that motivate the next contribution of this thesis.
The first limitation is that the code responsible for checking requirements must be directly inserted into the source code of the implementation under test.
This involves adding calls to a shared library at specific locations in the implementation.
Since these locations differ between implementations and between versions, the process requires a detailed understanding of the code base and must be repeated whenever the implementation changes.
This reduces the portability of the technique and increases the manual effort required to apply it.

The second limitation concerns the reliance on completing a protocol interaction.
In the case of DTLS, this means that SE must reach either the successful completion of the handshake or a definite failure triggered by the invalid input.
As a result, SE may be forced to explore many paths that lie far beyond the part of the execution that is relevant to the requirement being tested.
This unnecessarily increases the symbolic search space and can significantly slow down the detection of violating inputs.

These two limitations motivate the need for a technique that avoids direct instrumentation of the implementation and that can detect requirement violations without relying on the completion of a full protocol exchange.
The next chapter introduces a monitor based approach that addresses both issues by moving requirement checking outside the implementation and by enabling earlier detection of violation points.